\documentclass[10pt]{article}

% Specify the margins and text size.
\setlength{\textwidth}{6.4in}
\setlength{\textheight}{9.5in}
\setlength{\oddsidemargin}{0pt}
\setlength{\evensidemargin}{0pt}
\setlength{\topmargin}{0pt}
\setlength{\hoffset}{.05in}
\setlength{\voffset}{-.8in}

\setlength{\parskip}{5pt}
\setlength{\parindent}{0pt}

% Load some fonts and symbol packages
\usepackage{latexsym}
\usepackage{pifont}       % contains 'star' symbol for counterinsurgency handbook title
\usepackage{yfonts} 
\usepackage{amsmath}
\usepackage{amsfonts}

\usepackage{graphicx}     % actually, this is loaded by pstricks
\usepackage[T1]{fontenc}
\usepackage{ifthen}
%\usepackage{pstricks,pst-grad,pst-text,pst-node,multido,pst-plot,calc,pst-3dplot}
\usepackage{pstricks,pst-grad,pst-func,pst-text,pst-node,multido,pst-plot,calc,pst-3dplot}
%\usepackage[all]{xy}
%\usepackage{animate}

% The hyperref package inserts links into the pdf.
\definecolor{MyLinkColor}{rgb}{.1,.2,1}
\definecolor{MyCiteColor}{rgb}{.1,1,.2}
\definecolor{MyURLColor}{rgb}{.4,.4,.4}
\usepackage[backref=true,pagebackref=false,hyperindex,colorlinks=true,
  linkcolor=MyLinkColor,urlcolor=MyURLColor]{hyperref}


% The tweaklist package is something I found on the web.  It provides a simple interface
% for making changes to spacing used in the itemize and enumerate environments.  Comment
% this out if you don't care to use tweaklist.
\usepackage{tweaklist}
\renewcommand{\itemhook}{\setlength{\parskip}{2pt}\setlength{\parsep}%
{1pt}\setlength{\topsep}{0pt}\setlength{\itemsep}{0pt}}

\newcommand{\U}{\underline{\hspace{5pt}}}

\usepackage{listings}
\newcounter{EX}\setcounter{EX}{1}
\newcommand{\EXERCISE}{\arabic{EX}.\stepcounter{EX} }

\begin{document}
\pagestyle{empty}
\lstset{language=R, showspaces=false, showstringspaces=false}

\href{http://www.su.edu}{\includegraphics[height=1.75cm]{sulogo.eps}}
\vspace{-1.69cm}

{\small \hfill
\begin{tabular}{cl}
& Math 207\\& Introduction to Statistics\\
\end{tabular}
}
\setlength{\baselineskip}{1.05\baselineskip}
\medskip
\medskip

\begin{center}
\textbf{\large   NHANESIII Data Set:  Report II}
\end{center}

\newcommand{\SUBX}{\smallskip\hspace{10pt}}
\newcommand{\BSK}{\vspace{.14in}}

In a prior assignment you were asked to
 review the NHANES code books then decide on four questions about the data that interest 
you and that you would like to explore.  
You were asked to include variables of each of the four 
types (continuous, discrete, ordered and unordered). 

For this assignment, please prepare an RMarkdown file that answers the questions below.
Please include your R code.  Include text to explain what you are calculating and
what the results mean.  Recall that the following commands read the NHANES data files
into \texttt{R}.
\vspace{-6pt}
\begin{verbatim}
adult = read.csv("adult.csv", header=TRUE)
exam  = read.csv("exam.csv",  header=TRUE)
lab   = read.csv("lab.csv",   header=TRUE)
youth = read.csv("youth.csv", header=TRUE)
\end{verbatim}\vspace{-6pt}
Your operating system may require you to include full or partial path names rather
than just the file names.
\medskip

\EXERCISE \textbf{Normal Curve}. The standard normal (Gaussian) distribution is an ideal histogram which we can often use
as a model data.

\begin{center}
\begin{pspicture}(-3,-2.1)(3,1.8)
\psset{xunit=1.5,yunit=4,linewidth=0.02}
\psline(-3.5,0)(3.5,0)  
  \psline(-3,0)(-3,-0.04) \rput[t](-3,-0.07){\small -3\hphantom{-}}
  \psline(-2,0)(-2,-0.04) \rput[t](-2,-0.07){\small -2\hphantom{-}}
  \psline(-1,0)(-1,-0.04) \rput[t](-1,-0.07){\small -1\hphantom{-}}
  \psline(0,0)(0,-0.04)   \rput[t](0,-0.07){\small 0}
  \psline(1,0)(1,-0.04)   \rput[t](1,-0.07){\small 1}
  \psline(2,0)(2,-0.04)   \rput[t](2,-0.07){\small 2}
  \psline(3,0)(3,-0.04)   \rput[t](3,-0.07){\small 3}
  \rput[l](3.6,0){\small $x$}
\psline(0,0)(0,0.5)
  \psline(-0.12,0.5)(0,0.5)    \rput[r](-0.21,0.5){\small $0.5$}
  \psline(-0.12,0.25)(0,0.25)  \rput[r](-0.21,0.25){\small $0.25$}
\psGauss[linecolor=blue,linewidth=0.02,sigma=1,mue=0]{-3}{3}
\pnode(-1,-0.15){A}\pnode(1,-0.15){B}
\psbrace[braceWidth=0.02,braceWidthInner=5pt,braceWidthOuter=4pt](A)(B){\rput{90}(0.25,-0.05){\small 68\%}}
%
\pnode(-2,-0.15){C}\pnode(2,-0.15){D}
\psbrace[braceWidth=0.02,braceWidthInner=27pt,braceWidthOuter=4pt](C)(D){\rput{90}(0.25,-0.05){\small 95\%}}
%
\pnode(-3,-0.15){E}\pnode(3,-0.15){F}
\psbrace[braceWidth=0.02,braceWidthInner=50pt,braceWidthOuter=4pt](E)(F){\rput{90}(0.25,-0.1){\small 99.7\%}}
\end{pspicture}
\end{center}
\bigskip
\medskip


\hspace{10pt} a) Does the continuous (quantitative) variable that you chose to study
fit a normal curve?

\hspace{20pt} a.i) As a first test, compute the fraction of the data that is
within 1, 2 and 3 standard deviations of the mean. For example, \vspace{-6pt}
\begin{verbatim}
e1 = exam[exam$BMPBMI != 8888, ]
e1bmi = e1$BMPBMI
m  = mean(e1bmi, na.rm=T)
s  = sd(e1bmi, na.rm=T)
length(e1bmi[(m - s < e1bmi) & (e1bmi < m + s)]) / length(e1bmi)
# Repeat for 2*s and 3*s
\end{verbatim}
Do the percents roughly fit 68\%, 95\% and 99.7\%?  If so, what does this tell you about how
well the data fits a normal curve?
\vfill
\eject

\hspace{20pt} a.ii) As a second test, make a QQ plot comparing the quantiles of your data
to the quantiles of a normal curve.\vspace{-6pt}
\begin{verbatim}
qqnorm(e1bmi)
\end{verbatim}
If your data fits a normal curve, the QQ plot should be a straight line.  Comment on your plot
and the conclusion you reached in a.i).

\hspace{20pt} a.iii) The Kolmogorov-Smirnov test can be used to test the null hypothesis that
a given data set was generated from a specified distribution.  A larger $p$-value for the test
(e.g., at least 5\%) indicate that the data probably differs from the expected distribution
simply due to chance.  A small $p$-value suggests that the difference is not due to chance
and that the data probably came from some other distribution.
Run the test on your data:\vspace{-6pt}
\begin{verbatim}
ks.test(e1bmi, "pnorm", m, s)
\end{verbatim}
Report the $p$-value and comment on its significance.  How does your result relate to your
conclusions in a.i) and a.ii)?
\medskip

\hspace{10pt} b) Repeat part a) using the weight distribution variable
\texttt{exam\$BMPWT}.
\medskip
\bigskip

2. \textbf{Correlation}. 
In many situations, we may seek to use measurements of one or more variables to predict
the value of another variable.  For example, we could seek to predict the sale price of a 
house based on the sale prices of other houses in the neighborhood, size of the house, 
number of bedrooms, and interest rates.  We could predict the stopping distance of a car
using speed, weight, and road conditions.  Standardized test scores, high school grade point
average, and enrollment in advanced classes are used in college admissions because these
variables can be used to estimate college success.

The correlation coefficient $r$ is a measure of linear association between two lists
$\{x_1, \dots, \,x_n\}$ and $\{y_1, \dots,\,y_n\}$ of variables. It is defined by 
\begin{align*}
r &=\frac{1}{n-1}\,\sum_{i=1}^n\,\left(\frac{x_i - \mbox{mean}_x}{\mbox{sd}_x}\right)\,
  \left(\frac{y_i - \mbox{mean}_y}{\mbox{sd}_y}\right) =
\mbox{mean of the $x$ and $y$ values measured in standard units}%\vspace{-20pt}
\end{align*}\vspace{-15pt}

Observations about $r$:
\begin{itemize}
\item The correlation $r$ is a number between $-1$ and $1$.
%\item It is the average product of $z$-scores for the $X$ and $Y$ variables.
\item It measures linear association between the lists.
\item It measures the clustering of the $(x_i, y_i)$ points around a line.
\item If $r$ is close to $1$ or $-1$, the points fit rather tightly along a line.
%\item For values of $r$ near $0$, the points form more of a cloud with (perhaps) and up
%   or down trend.
%\item A negative correlation means that the cloud of $(x_i, y_i)$ points slopes down.
\item The correlation between $X$ and $Y$ is equal to the correlation between $Y$ and $X$.
\end{itemize}
\smallskip

\hspace{10pt} a) Make a plot of your two quantitative variables.  Be sure to label the axes.
\vspace{-6pt}
\begin{verbatim}
plot( x variable, y variable, xlab="X axis label", ylab="Y axis label")
\end{verbatim}
Does there appear to be a linear association between the variables?  
If so, is it strong (close to $\pm 1$) or weak (closer to 0)?
\medskip

\hspace{10pt} b) Compute the correlation between the two variables and comment on your 
findings. What, if anything, do the sign and magnitude of the correlation indicate?\vspace{-6pt}
\begin{verbatim}
cor(x variable, y variable)
cor(y variable, x variable)
\end{verbatim}

\hspace{10pt} c)  A strong correlation can indicate that one variable can cause the 
other (e.g., between hours spent reading and student GPA).  It 
can mean that there is some underlying common cause for both (e.g., between drinking and lung
cancer where the confounding factor is an association between smoking and drinking).  
It can also be simply a statistical artifact (e.g., between sunspot activity and stock returns
over some conveniently selected time periods).  Does the correlation between your variables
suggest a cause and effect relationship, an underlying confounding factor or a statistical
artifact?
\vfill
\eject


\end{document}

